\subsection{DEFINITION OF SOLVABLE LIE ALGEBRAS}

DEFINITION 1. A Lie algebra g is called solvable if its kth derived algebra \(\mathcal{D}^{k}\mathfrak{g}\) is zero for sufficiently large \(k\) .

A nilpotent Lie algebra is solvable.

PROPOSITION 1. Subalgebras and quotient algebras of a solvable Lie algebra are solvable. Every extension of a solvable algebra by a solvable algebra is solvable. Every finite product of solvable algebras is solvable.

Let g be a Lie algebra, \(g'\) a subalgebra, h an ideal of g, t = g/h and \(\phi\) the canonical mapping of g onto t. If g is solvable then \(D^{k}g = \{0\}\) for some integer k, hence \(D^{k}g' \subset D^{k}g = \{0\}\) , \(D^{k}t = \phi(\mathcal{D}^{k}g) = \{0\}\) and hence \(g'\) and t are solvable. If h and t are solvable, there exist integers s, t such that

\[
\mathcal {D} ^ {s} \mathfrak {h} = \mathcal {D} ^ {t} \mathfrak {k} = \{0 \};
\]

then \(\mathcal{D}^t\mathfrak{g}\subset \mathfrak{h}\) , hence \(\mathcal{D}^{s + t}\mathfrak{g} = \mathcal{D}^{s}(\mathcal{D}^{t}\mathfrak{g})\subset \mathcal{D}^{s}\mathfrak{h} = \{0\}\) and \(\mathfrak{g}\) is solvable. The last assertion follows from the second by induction on the number of factors.

PROPOSITION 2. Let g be a Lie algebra. The following conditions are equivalent:

\begin{enumerate}
\def\labelenumi{(\alph{enumi})}
\item
  g is solvable;
\item
  there exists a decreasing sequence \(\mathfrak{g} = \mathfrak{g}_0 \supset \mathfrak{g}_1 \supset \dots \supset \mathfrak{g}_n = \{0\}\) of ideals of \(\mathfrak{g}\) such that the algebras \(\mathfrak{g}_i / \mathfrak{g}_{i+1}\) are commutative ( \(i = 0, 1, \ldots, n-1\) );
\item
  there exists a decreasing sequence \(\mathfrak{g} = \mathfrak{g}_0' \supset \mathfrak{g}_1' \supset \dots \supset \mathfrak{g}_p' = \{0\}\) of subalgebras of \(\mathfrak{g}\) such that \(\mathfrak{g}_{i+1}'\) is an ideal of \(\mathfrak{g}_i'\) and \(\mathfrak{g}_i'/\mathfrak{g}_{i+1}'\) is commutative ( \(i = 0, 1, \ldots, p - 1\) );
\item
  there exists a decreasing sequence \(\mathfrak{g} = \mathfrak{g}_0'' \supset \mathfrak{g}_1'' \supset \dots \supset \mathfrak{g}_q'' = \{0\}\) of subalgebras of \(\mathfrak{g}\) such that \(\mathfrak{g}_{i+1}''\) is an ideal of \(\mathfrak{g}_i''\) of codimension 1 ( \(i = 0, 1, \ldots, q - 1\) ).
\item
  \(\Rightarrow\) (b): it suffices to consider the sequence of derived ideals of g.
\item
  \(\Rightarrow\) (c): this is obvious.
\item
  \(\Rightarrow\) (d): suppose that condition (c) holds; every vector subspace of \(\mathfrak{g}_i'\) containing \(\mathfrak{g}_{i+1}'\) is an ideal of \(\mathfrak{g}_i'\) , whence immediately (d).
\item
  \(\Rightarrow\) (a): this follows immediately from the fact that an extension of a solvable algebra by a solvable algebra is solvable.
\end{enumerate}

